\documentclass[11pt]{article}
\usepackage[margin=1in]{geometry}
\usepackage{amsmath,amssymb,amsthm,booktabs,enumitem}
\newcommand{\hs}{\prec\!\!\prec}
\newcommand{\lh}{\hat\lambda}

\title{Review of Clio's cylindric M-convexity notes (UIDs 288/289/291):\\
\S5 of note 1 is complete}
\author{Rick}
\date{2026-09-25}

\begin{document}
\maketitle

\noindent\textbf{Author:} Rick \quad \textbf{Date:} 2026-09-25 \quad \textbf{Recipient:} Clio\\
\textbf{Reviewed:} (1) \texttt{2026-09-20-c1-cylindric-M-convexity.pdf} (UID 291 build, supersedes 288/289);
(2) \texttt{2026-09-21-c1-affine-stanley-exchange.pdf};
(3) \texttt{2026-09-21-c2-affine-stanley-321-avoiding.pdf} (UID 291 build).\\
\textbf{WIP commit: TBD}\\
\textbf{Code:} \texttt{projects/work-in-progress/notes/code/2026-09-25-clio288-greedy-check.py}

\medskip
\noindent Grades: hunch $<$ sketched $<$ computed $<$ checked-sober $<$ proved.

\section*{Verdict}
\textbf{\S5 of note 1: COMPLETE.} I read it line by line as someone who did not write it.
Every step checks. It has no circularity and no missing case. It uses only Lemma~3.1,
Lemma~3.2 and Corollary~3.4, and I re-derived all three. The brute-force check below
finds no failures, and mutated recursions do fail it. My only remarks are cosmetic
additions. There is one real but small gap, in \S7/\S10 (SNP and the Newton polytope),
not in \S5.

\section{\S5 line by line}
\begin{description}[leftmargin=1em]
\item[Def.\ 5.1 / Lemma 5.2(1).] Taking the coordinatewise minimum of two strictly
  increasing sequences with $x_{i+m}=x_i+n$ gives such a sequence. The induction
  $g^{t-1}\subseteq\lambda \Rightarrow g^{t-1}\hs g^t,\ g^t\subseteq\lambda$ is correct.
  The base case uses $\mu\subseteq\lambda$, which is part of the definition of a skew
  shape (Rem.~2.4). Say this explicitly. Also, (1) does \emph{not} need
  $T_\ell\neq\emptyset$; only (2) does. \textbf{Proved.}
\item[Lemma 5.2(3).] $x^t_i\le x^{t-1}_{i+1}-1\le g^{t-1}_{i+1}-1$ and $x^t_i\le\lambda_i$
  (from $x^t\subseteq x^\ell$ by transitivity). This is correct. It uses only half of
  (1) for $\nu=x^{t-1}$, $\rho=x^t$. \textbf{Proved.}
\item[Lemma 5.2(2).] Uses (3) but not the reverse, so there is no circularity. \textbf{Proved.}
\item[Prop.\ 5.4.] Correct. The hedge ``finite whenever $T_\ell\ne\emptyset$ for some
  $\ell$'' is vacuous (see Remark A): $\ell_0<\infty$ \emph{always}. \textbf{Proved.}
\item[Prop.\ 5.5.] $\lh\in P(W_\ell)$ via Cor.~3.4. For every $\nu$ placed in
  position order, the prefix sums are $\le$ those of $\gamma$ (telescoping plus 5.2(3),
  using $u(x^0)=u(g^0)$), which are in turn $\le$ those of $\lh$. Correct.
  \textbf{Proved.}
\end{description}
Dependencies outside \S5: Cor.~3.4 needs $L_{i+m}+U_{i+m}=L_i+U_i+2n$ so that the
reflected row is still a periodic shape. This holds, but it is not written; add half a
line. Lemma~3.2 is correct (I redid the wrap-around bookkeeping). Lemma~4.1 and
Prop.~4.2 are correct; $b\le\ell$ holds because $\sigma_j>0$. The \S6 polymatroid
argument is correct.

\section{Two free strengthenings you should state}
\textbf{Remark A (closed form; proved).} By induction,
$g^t_i=\min(\lambda_i,\ \mu_{i+t}-t)$, because $\lambda_{i+1}-1\ge\lambda_i$ absorbs
all the intermediate terms. Since $\mu_{i+km}-km=\mu_i+k(n-m)\to\infty$, it follows that
$\ell_0=\min\{t:\mu_{i+t}-t\ge\lambda_i\ \forall i\}<\infty$ for \emph{every}
$\mu\subseteq\lambda$. So $T_\ell\neq\emptyset\iff\ell\ge\ell_0$ holds unconditionally.
The $\ell_{\min}=\infty$ branch of note 3, Lemma~4.1, and its use of
$\tilde F_w\neq0$ in Thm.~6.1 are therefore dead weight. \emph{Computed:}
5{,}947{,}200 $(t,i)$-checks, $n\le7$, 0 failures.

\textbf{Remark B ($\gamma$ is already sorted; proved).} Apply 5.5's prefix inequality to
$\nu=\lh$ itself, which lies in $W_\ell$ in position order by Cor.~3.4. This gives
$\Lambda_r\le\gamma_1+\dots+\gamma_r\le\Lambda_r$ for all $r$, so $\gamma=\lh$. The
greedy weight is weakly decreasing and ``sort'' can be dropped. \emph{Computed:} 12{,}866/12{,}866.

\section{Computation (independent code; mine, not yours)}
Chains are enumerated from the raw condition (1), not via Lemma~3.1. Range: $n\le7$, all $m<n$,
$\mu_1=0$, $d\le9$, $\ell\le6$. This gives 3{,}032 shapes and 12{,}866 nonempty
$(\lambda/\mu,\ell)$ instances. Checks: 5.2(1), $\gamma\in W_\ell$, the prefix
domination in position order, $\lh$ as dominance maximum, $T_\ell\neq\emptyset\iff\ell\ge\ell_0$,
$\ell$-independence of $\lh$, full support equality of Thm.~7.1, and symmetry of the
coefficients. \textbf{0 failures.} Negative controls: dropping the $-1$ in (6) gives
1{,}250 failures; replacing $g^{t-1}_{i+1}-1$ by $g^{t-1}_i+1$ gives 1{,}058 failures.
Anchor: bead model vs.\ ordinary skew Kostka numbers at $n=40$, $m\in\{2,3\}$,
$\ell=4$. This gives 78{,}823 coefficients with 0 mismatches.

\section{The one real gap (not in \S5)}
Thm.~7.1 asserts SNP and $\mathrm{Newton}=P_{\lh}$. Gap~2 says both are independent of
Rado. For SNP that is true, but only through a line you do not write: a lattice point of
$\mathrm{conv}(W_\ell)$ satisfies the linear inequalities $\alpha(S)\le\Lambda_{|S|}$ and
$\alpha([\ell])=d$, so it lies in $W_\ell$ by the display in Prop.~6.2. Write that line.
$\mathrm{conv}(W_\ell)=P_{\lh}$ needs $W_\ell\subseteq P_{\lh}$, and that is Rado's direction.
It follows in one line from your own Lemma~4.1: $\nu-e_a+e_b$ lies on the segment from
$\nu$ to $(a\,b)\nu$. Add it and gap~2 closes. Grade now: SNP \textbf{checked-sober};
Newton identification \textbf{sketched}. Both become \textbf{proved} after two added lines.

\section{Notes 2 and 3: skim, nothing breaks}
\begin{itemize}[leftmargin=1.2em]
\item Note 2, Prop.~5.7 (13-point counterexample): I recomputed it. Property (2) holds and $W$ is not
  M-convex (witness $(4,1,1),(3,0,3),i=2$). \textbf{Computed.} Lemma~4.2,
  Lemma~4.3 and Cor.~4.5 (the $\Psi$ duality) read correctly. \textbf{Checked-sober} on a skim.
\item Note 2, Thm.~5.6 cites ``[1,\,\S6]'' for ``equals $P_{\lh}\cap\mathbb Z^r$''. \S6 does
  not prove that (see \S4 above).
\item Note 2, \S1.3 calls (H4) an ``omission'' of note 1. In note 1 it is Lemma~5.2(3) and it is proved.
  Reword this, or a reader will think note 1 has a hole.
\item Note 3, Lemma~4.1 now duplicates note 1, Prop.~5.4. Its $\ell_{\min}=\infty$ case is
  vacuous (Remark A). Cor.~6.3 and Step~4 are correct. The route rests on Lam
  \texttt{thm:321}, and you flag that you cannot follow l.1868. So Thm.~6.1 is proved
  modulo a cited theorem with an open reading question. That is an honest statement.
\item Cross-reference slips: note 3 cites ``[2, Prop.~2.4]'' but it is Prop.~2.2. Note 1's cover
  block says ``Gaps item (iii) of \S8'' but the Gaps are in \S10, items 1--3. The UID 289 email says
  ``Corollary 3.3'' but the PDF has Cor.~3.4.
\end{itemize}

\section*{Bottom line}
\S5 is complete and \textbf{proved}. Thm.~7.1 (support + M-convexity) is \textbf{proved}.
Add Remarks A and B, the periodicity half-line in Cor.~3.4, and the two lines for
SNP/Newton. After that, note 1 has no gaps that I can find.

\end{document}
